\documentclass{beamer}
\usepackage{lmodern}

\usepackage{xeCJK}

\usepackage[safe]{silence}
\WarningFilter*{latexfont}{Font}

\usetheme{CambridgeUS} % try Madrid, Pittsburgh
\usecolortheme{beaver}
\usefonttheme[onlymath]{serif} % try "professionalfonts"

\setbeamertemplate{theorems}[normal font]
\setbeamertemplate{itemize items}[default]
\setbeamertemplate{enumerate items}[default]

\usepackage{amsmath, amsfonts, latexsym, mathtools, extarrows}
% \usepackage{kbordermatrix}
% \usepackage{ntheorem}

\usepackage{centernot}

\DeclareMathOperator*{\argmin}{arg\,min}
\DeclareMathOperator*{\argmax}{arg\,max}

% colors
\newcommand{\red}[1]{\textcolor{red}{#1}}
\newcommand{\green}[1]{\textcolor{green}{#1}}
\newcommand{\blue}[1]{\textcolor{blue}{#1}}
\newcommand{\purple}[1]{\textcolor{purple}{#1}}

% color box
\newcommand{\rbox}[1]{\red{\boxed{#1}}}
\newcommand{\gbox}[1]{\green{\boxed{#1}}}
\newcommand{\bbox}[1]{\blue{\boxed{#1}}}
\newcommand{\pbox}[1]{\purple{\boxed{#1}}}

\usepackage{tikz}
% see http://tex.stackexchange.com/a/7045/23098
\newcommand*\circled[1]{\tikz[baseline=(char.base)]{
            \node[shape=circle,draw,inner sep=2pt] (char) {#1};}}
\usepackage{tikz-qtree}
\usetikzlibrary{backgrounds, fit}

\usepackage{pifont}
\usepackage{wasysym}

\usepackage[normalem]{ulem}
\newcommand{\middlewave}[1]{\raisebox{0.5em}{\uwave{\hspace{#1}}}}

\usepackage{graphicx, subcaption}

\usepackage{algorithm}
\usepackage[noend]{algpseudocode}

\newcommand{\pno}[1]{\textcolor{blue}{\scriptsize [Problem: #1]}}
\newcommand{\set}[1]{\{#1\}}
\DeclareMathOperator{\lcm}{lcm}

\newcommand{\cmark}{\green{\ding{51}}}
\newcommand{\xmark}{\red{\ding{55}}}
%%%%%%%%%%%%%%%%%%%%%%%%%%%%%%%%%%%%%%%%%%%%%%%%%%%%%%%%%%%%%%
% for fig without caption: #1: width/size; #2: fig file
\newcommand{\fignocaption}[2]{
  \begin{figure}[htp]
    \centering
      \includegraphics[#1]{#2}
  \end{figure}
}

% for fig with caption: #1: width/size; #2: fig file; #3: fig caption
\newcommand{\fig}[3]{
  \begin{figure}[htp]
    \centering
      \includegraphics[#1]{#2}
      \caption[labelInTOC]{#3}
  \end{figure}
}

\newcommand{\p}{\text{P}}
\newcommand{\np}{\text{NP}}
\newcommand{\conp}{\text{coNP}}
\newcommand{\pr}{\le_{p}} % polynomial reduction
\newcommand{\nph}{\text{NP-hard}}
\newcommand{\npc}{\text{NP-complete}}
\newcommand{\npi}{\text{NP-intermediate}}
\newcommand{\conpc}{\text{coNP-complete}}
\newcommand{\expcls}{\text{EXP}}
\newcommand{\rcls}{\text{R}}
\newcommand{\dtime}{\text{DTIME}}
\newcommand{\dspace}{\text{SPACE}}
\newcommand{\pspace}{\text{PSPACE}}
\newcommand{\pspacec}{\text{PSPACE-complete}}
\newcommand{\ph}{\text{PH}}
\newcommand{\hp}{\text{HAM-PATH}}
\newcommand{\hc}{\text{HAM-CYCLE}}
\newcommand{\fp}{\text{FP}}
\newcommand{\poly}{\text{Poly}}
\newcommand{\po}{\text{PO}}
\newcommand{\fptas}{\text{FPTAS}}
\newcommand{\ptas}{\text{PTAS}}
\newcommand{\apx}{\text{APX}}
\newcommand{\logapx}{\text{Log-APX}}
\newcommand{\polyapx}{\text{Poly-APX}}
\newcommand{\expapx}{\text{Exp-APX}}
\newcommand{\npo}{\text{NPO}}
\newcommand{\ms}{\text{MS}}

\newcommand{\titletext}{Approximation Algorithms}

\newcommand{\tx}{%
  \section*{}
  \begin{frame}[noframenumbering]
	\fignocaption{width = 0.50\textwidth}{figs/thankyou.jpg}
  \end{frame}
}
%%%%%%%%%%%%%%%%%%%%
\title[\titletext]{\titletext}
% \subtitle{}

\author[Hengfeng Wei]{Hengfeng Wei}
\titlegraphic{\includegraphics[height = 1.5cm]{figs/qrcode-approx-alg-20170518.png}}
\institute{hfwei@nju.edu.cn}
\date{May 14 $\sim$ May 18, 2017}

\AtBeginSection[]{
  \begin{frame}[noframenumbering, plain]
    \frametitle{\titletext}
    \tableofcontents[currentsection, sectionstyle=show/shaded, subsectionstyle=show/show/hide]
  \end{frame}
}
%%%%%%%%%%
\begin{document}
\maketitle

\section{Approximation Classes}

%%%%%%%%%%%%%%%%%%%%
\begin{frame}{Approximation Classes}
  \begin{description}
	\item[NPO:] \np{} optimization problems
	\item[APX:] constant factor $\epsilon$-approximation ($\epsilon > 1$)
	\item[$f(n)$-APX:] Exp-APX, Poly-APX, Log-APX
	\item[PTAS:] 
	  \begin{itemize}
		\item $\forall \epsilon > 0: (1 + \epsilon)$-approximation
		\item $\p: \poly(n) \qquad O(n^{2/\epsilon}) \qquad O(n^{2^{2^{1/\epsilon}}})$
	  \end{itemize}
	\item[FPTAS:] 
	  \begin{itemize}
		\item $\forall \epsilon > 0: (1 + \epsilon)$-approximation
		\item $\fp: \poly(n, 1/\epsilon) \qquad O((1/\epsilon)^2 \cdot n^3)$
	  \end{itemize}
	\item[PO:] polynomial time solvable optimization problems
  \end{description}
\end{frame}
%%%%%%%%%%%%%%%%%%%%
\begin{frame}{Approximation Classes}
  \begin{align*}
	\po &\subsetneq \fptas \subsetneq \ptas \subsetneq \apx \\
	&\subsetneq \logapx \subsetneq \polyapx \subsetneq \expapx \subsetneq \npo
  \end{align*}

  \vspace{-0.50cm}
  \[
	(\text{if } \p{} \neq \np{})
  \]

  \begin{enumerate}
	\item $\text{Knapsack} \in \fptas \setminus \po$
	\item $\text{Makespan} \in \ptas \setminus \fptas$ ({\small \textcolor{blue}{TODAY}})
	\item $\text{Vertex Cover} \in \apx \setminus \ptas$ 
	\item $\text{Set Cover} \in \logapx \setminus \apx$ ({\small CLRS 35.3})
	\item $\text{Clique} \in \polyapx \setminus \logapx$
	\item $\text{TSP} \in \expapx \setminus \polyapx$
  \end{enumerate}

  \begin{alertblock}{Reference}
	\begin{itemize}
	  \item ``A Survey on the Structure of Approximation Classes'' by Bruno Escoffier, 	Vangelis Th. Paschos, 2010.
	\end{itemize}
  \end{alertblock}
\end{frame}
%%%%%%%%%%%%%%%%%%%%
\begin{frame}{Stability of approximation}
  TSP: worst-case complexity \emph{vs.} inapproximability according to instances

  \begin{itemize}
	\item $\text{TSP} \in \expapx \setminus \polyapx$
	\item $\Delta\text{-TSP} \in \apx{}$
	\item $\text{Euclidean TSP} \in \ptas{}$
  \end{itemize}

  \begin{alertblock}{Reference}
	\begin{itemize}
	  \item ``Stability of Approximation Algorithms for Hard Optimization Problems'' by Juraj Hromkovi\v{c}, 1999.
	\end{itemize}
  \end{alertblock}
\end{frame}
%%%%%%%%%%%%%%%%%%%%
\begin{frame}{Stability of approximation}
  \begin{exampleblock}{Distance function (JH 4.2.3.3)}
	\begin{gather*}
	  \text{dist}_{k}(G,c) = \\
	  \max\left\{0, \max\left\{ \frac{c(\set{u,v})}{\sum_{i=1}^{m} c(\set{p_i, p_{i+1}})} - 1 \bigl \lvert |u = p_1 \leadsto v = p_m| \le k \right\}\right\}
	\end{gather*}
  \end{exampleblock}

  \begin{center}
	\sout{enumerate}: $k = n^{\frac{1}{3}}$ \\[5pt]
    shortest paths of length $\le k$ (Bellman-Ford)
  \end{center}
\end{frame}
%%%%%%%%%%%%%%%%%%%%
\begin{frame}{Stability of approximation}
  \begin{exampleblock}{$h_{\text{index}}$ (JH 4.2.3.4)}
	\[
	  h_{\text{index}}\text{: using canonical order}
	\]
  \end{exampleblock}

  \[
	|\text{Ball}_{r, h_{\text{index}}}(L_I)| < \infty
  \]

  \[
	\delta_{r, \epsilon} = \max \set{R_A(x): x \in \text{Ball}_{r, h_{\text{index}}}(L_I)}
  \]
\end{frame}
%%%%%%%%%%%%%%%%%%%%
\begin{frame}{Stability of approximation}
  \begin{exampleblock}{$h$ (JH 4.2.3.5)}
	\begin{itemize}
	  \item $h$: infinite jumps
	  \item $\delta$-approx. algorithm $A$ for $U$ is stable according to $h$
	\end{itemize}
  \end{exampleblock}

  \[
	\text{TSP } U:  (G, 1)
  \]

  \[
	\text{Multi-TSP } \overline{U}: (G, k)
  \]

  \[
	h(G,k) = k-1
  \]

  \[
	A \text{ is } \delta\text{-approx. for } (G, 1) \implies A \text{ is } r\delta\text{-approx. for } (G,r \in \mathbb{N})
  \]
\end{frame}
%%%%%%%%%%%%%%%%%%%%

\section{Approximation Algorithms}

%%%%%%%%%%%%%%%%%%%%%%%%%%%%%%
% \subsection{Makespan Scheduling Problem}

%%%%%%%%%%%%%%%%%%%%
\begin{frame}{Makespan scheduling problem}
  \begin{exampleblock}{Makespan scheduling problem (\ms{})}
	\begin{itemize}
	  \item $n$ jobs: $J_1, \dots, J_n$
	  \item processing time: $p_1, \dots, p_n$
	  \item $m \ge 2$ machines: $M_1, \dots, M_m$
	 ­\item goal: minimize the makespan
	\end{itemize}
  \end{exampleblock}
\end{frame}
%%%%%%%%%%%%%%%%%%%%
\begin{frame}{\ms{} is \npc{}}
  \begin{definition}[Partition]
	\begin{description}
	  \item[Instance:] 
		\[
		  \forall a \in A: s(a) \in \mathbb{Z}^{+} 
		\]
	  \item[Question:] Is there a subset $A' \subseteq A$: 
		\[
		  \sum_{a \in A'} s(a) = \sum_{a \in A \setminus A'} s(a)
		\]
	\end{description}
  \end{definition}

  \[
	A = \set{5,1,3,4,8,2,7} \implies A' = \set{5,3,7}, A \setminus A' = \set{1,2,4,8}
  \]
\end{frame}
%%%%%%%%%%%%%%%%%%%%
\begin{frame}{\ms{} is strongly \npc{}}
  \begin{definition}[$3$-Partition]
	\begin{description}
	  \item[Instance:] 
		\begin{gather*}
		  |A| = 3m, B \in \mathbb{Z}^{+} \\
		  \forall a \in A: s(a) \in \mathbb{Z}^{+}, B/4 < s(a) < B/2
		\end{gather*}
	  \item[Question:] Can $A$ be partitioned into $m$ disjoint sets $S_1,\dots,S_m$:
		\[
		  \forall 1 \le i \le m: |S_i| = 3, \sum_{a \in S_i} s(a) = B
		\]
	\end{description}
  \end{definition}

  \[
	A = \set{1,2,2,3,3,4,6,7,8}, m = 3, B = 12 \implies \set{1,3,8; 2,4,6; 2,3,7}
  \]
\end{frame}
%%%%%%%%%%%%%%%%%%%%
\begin{frame}{List-Scheduling (LS) algorithm}
  \begin{exampleblock}{List-Scheduling algorithm (JH 4.2.1.4)}
	\begin{itemize}
	  \item online
	  \item assign job to the least heavily loaded
	\end{itemize}
  \end{exampleblock}

  \[
	J = {2,3,4,6,2,2} \qquad m = 3
  \]
\end{frame}
%%%%%%%%%%%%%%%%%%%%
\begin{frame}{LS is $2$-approx.}
  \[
	T \emph{ vs. } T^{\ast}: \frac{T}{T^{\ast}}
  \]

  \begin{align*}
	T^{\ast} &\ge \frac{1}{m} \sum_{j} t_j \\
	T^{\ast} &\ge \max_{j} t_j
  \end{align*}

  \begin{enumerate}
	\item $M_i$: with the maximum load
	\item $J_i$: the last job on $M_i$
  \end{enumerate}

  \[
	ms_i \le \sum_{j} t_j \implies s_i \le \frac{1}{m} \sum_{j} t_j \le T^{\ast}
  \]

  \[
	T = c_i = s_i + p_i \le T^{\ast} + T^{\ast} = 2T^{\ast}
  \]
\end{frame}
%%%%%%%%%%%%%%%%%%%%
\begin{frame}{$2$-approx. is (almost) tight}
  \[
	n = \underbrace{m(m - 1)}_{p_i = 1} + \underbrace{1}_{p_i = m}
  \]

  \[
	\frac{T}{T^{\ast}} = \frac{2m - 1}{m} = 2 - \frac{1}{m}
  \]
\end{frame}
%%%%%%%%%%%%%%%%%%%%
\begin{frame}{LS is $(2-\frac{1}{m})$-approx.}
  \begin{align*}
	m s_i &\le \sum_{j \neq i} p_j = \frac{1}{m} (\sum_{j} p_j - p_i) \\
	&= \frac{1}{m} \sum_{j} p_j - \frac{1}{m} p_i \\
	&\le T^{\ast} - \frac{1}{m} p_i
  \end{align*}

  \begin{align*}
	T &= c_i = s_i + p_i \\
	&\le T^{\ast} + (1 - \frac{1}{m}) p_i
  \end{align*}
\end{frame}
%%%%%%%%%%%%%%%%%%%%
\begin{frame}{Sorting-Scheduling algorithm}
  \begin{exampleblock}{Sorting-Scheduling algorithm (JH 4.2.1.5)}
	Longest Processing Time (LPT) rule:
	\begin{itemize}
	  \item sorting non-increasingly
	  \item applying LS
	\end{itemize}
  \end{exampleblock}

  \begin{enumerate}
	\item $M_i$: with the maximum load
	\item $J_i$: the last job on $M_i$
  \end{enumerate}

  \[
	|M_i| = 1 \implies T = T^{\ast}
  \]

  \begin{align*}
	|M_i| \ge 2 &\implies p_i \le \frac{1}{2} T^{\ast} \\
	& \implies T = s_i + p_i \le (\frac{3}{2} - \frac{1}{2m}) T^{\ast}
  \end{align*}
\end{frame}
%%%%%%%%%%%%%%%%%%%%
\begin{frame}{LPT rule is $(\frac{4}{3} - \frac{1}{3m})$-approx.}
  \[
	p_1 \ge p_2 \ge \cdots \ge p_n
  \]

  \begin{description}
	\item[CASE $p_i \le \frac{1}{3} T^{\ast}$:] 
	  \[
		T \le (\frac{4}{3} - \frac{1}{3m}) T^{\ast}
	  \]
	\item[CASE $p_i > \frac{1}{3} T^{\ast}$:]
	  \begin{align*}
		p_i &\equiv p_n \text{\small (\textcolor{red}{w.l.o.g}: $T$ unchanged; $T^{\ast}$ not smaller)} \\
		&\implies p_1 \ge p_2 \ge \cdots \ge p_n > \frac{1}{3} T^{\ast} \implies |M_i| \le 2 \\
		&\implies n \le 2m \implies n = 2m - h \xLongrightarrow[\text{\small argument}]{\text{\small exchange}} T = T^{\ast}
	  \end{align*}
  \end{description}
\end{frame}
%%%%%%%%%%%%%%%%%%%%
\begin{frame}{$(\frac{4}{3} - \frac{1}{3m})$-approx. is tight}
  \[
	n = 2m + 1
  \]

  \begin{align*}
	\text{LPT: } &J_1 = x, J_{2m} = y, J_{2m+1} = z \\
	\text{OPT: } &J_{2m-1} = J_{2m} = J_{2m+1} = y \\
	\text{\it vs. } &\frac{x + 2y}{3y} = \frac{4}{3} - \frac{1}{3m}
  \end{align*}

  \[
	J = \set{2m-1,2m-1,\dots,m+1,m+1,m,m,m}
  \]

  \begin{alertblock}{Reference}
	\begin{itemize}
	  \item ``Bounds on Multiprocessing Timing Anomalies'' by R.L.Graham, 1969
	  \item ``Approximation Algorithms for NP-Hard Problems'' edited by Dorit Hochbaum, 1996 ({\small Theorem 1.5})
	\end{itemize}
  \end{alertblock}
\end{frame}
%%%%%%%%%%%%%%%%%%%%
\begin{frame}{\ptas{} for \ms{}}
  \begin{enumerate}
	\item $M_i$: with the maximum load
	\item $J_i$: the last job on $M_i$
  \end{enumerate}

  \[
	T = s_i + p_i \le T^{\ast} + p_i
  \]

  \begin{enumerate}
	\item $J = J_{L} \triangleq \set{\text{long jobs}} \uplus J_{S} \triangleq \set{\text{short jobs}}$
	\item $S_L$: the optimal schedule for $J_L$
	\item $S$: apply List-Scheduling to $S_L$ and $J_S$
  \end{enumerate}

  \begin{alertblock}{Reference}
	\begin{itemize}
	  \item ``The Design of Approximation Algorithms'' by David P. Williamson and David Shmoys, 2011 ({\small Section 3.2}).
	\end{itemize}
  \end{alertblock}
\end{frame}
%%%%%%%%%%%%%%%%%%%%
\begin{frame}{\ptas{} for \ms{}}
  \begin{enumerate}
	\item Split $J$:
	  \begin{gather*}
		J_i \in J_{S} \iff p_i \le \epsilon \cdot \frac{1}{m} \sum_{j} p_j \\
		\implies |J_L| < \frac{1}{\epsilon} \cdot m
	  \end{gather*}
	\item Time for $S_L$ (\textcolor{red}{$m$ being a constant!}):
	  \[
		m^{\frac{1}{\epsilon} \cdot m} \cdot O(n)
	  \]
	\item Approx. ratio ($p_i \in J_{S}$ case):
	  \begin{align*}
		T &= s_i + p_i \le \frac{1}{m} \sum_{j} p_j + \epsilon \cdot \frac{1}{m} \sum_{j} p_j \\
		&= (1 + \epsilon) \frac{1}{m} \sum_{j} p_j \\
		&\le (1 + \epsilon) T^{\ast}
	  \end{align*}
  \end{enumerate}
\end{frame}
%%%%%%%%%%%%%%%%%%%%
\begin{frame}{No \fptas{} for \ms{}}
  \begin{theorem}[$\ms \in \ptas \setminus \fptas$]
	No \fptas{} for \ms{}.
  \end{theorem}

  \[
	\ms{} \text{ is strongly } \npc{} \implies \ms{} \text{ with } \max_{j} p_j \le q(n) \text{ is } \npc{}.
  \]

  \begin{alertblock}{Reference}
	\begin{itemize}
	  \item ``The Design of Approximation Algorithms'' by David P. Williamson and David Shmoys, 2011 ({\small Section 3.2}).
	\end{itemize}
  \end{alertblock}
\end{frame}
%%%%%%%%%%%%%%%%%%%%
\begin{frame}{No \fptas{} for \ms{}}
  \begin{theorem}[]
	\[
	  \exists \fptas{} \text{ for } \ms{} \implies \ms \in \p{}.
	\]
  \end{theorem}

  \[
	A_{\epsilon}: \epsilon = \frac{1}{\lceil 2n q(n) \rceil}
  \]

  \begin{align*}
	(1 + \epsilon) T^{\ast} &= T^{\ast} + \epsilon \cdot T^{\ast} \\
	&\le T^{\ast} + \frac{1}{\lceil 2n q(n) \rceil} \cdot n q(n) \\
	&\le T^{\ast} + \frac{1}{2}
  \end{align*}

  \[
	\text{Time: } \text{Poly}(\frac{1}{\epsilon}, n) = \text{Poly}(\lceil 2n q(n) \rceil, n) = \text{Poly}(n)
  \]
\end{frame}
%%%%%%%%%%%%%%%%%%%%
%%%%%%%%%%%%%%%%%%%%%%%%%%%%%%

\tx{}

\end{document}
%%%%%%%%%%
